\section{Signed cubic certificates with exact quadratic moment information}
\label{sec:xor-quadratic}
The clique inequality in Proposition~\ref{prop:clique-escape} closes the
fractional-cardinality example. A signed cubic objective allows a different
obstruction: all first and second moments can be genuine moments while
the three-variable correlations still give a false objective value. The
Boolean moment construction used below is classical
\citep{Schoenebeck2008}. We prove its transfer to continuous coordinate
regions, with the exact quadratic moment hull imposed at each region.

Let $E$ be a multiset of $m>0$ clauses on $[n]$, each containing three
\emph{distinct} coordinates, with signs $b_e\in\{-1,1\}$. Occurrences count
multiplicity, and $\Delta=\max_i|\{e:i\in e\}|\ge1$. Set
\begin{equation}\label{eq:xor-objective}
 c_e(x)=\frac{1-b_e\prod_{i\in e}x_i}{2},\qquad
 F(x)=\frac1m\sum_{e\in E}c_e(x),\qquad x\in[-1,1]^n.
\end{equation}
Every $c_e$ lies in $[0,1]$ on the cube. Successively minimizing in each
coordinate moves to a Boolean vertex without increasing $F$, so
\begin{equation}\label{eq:xor-vertex-min}
 \OPT=\min_{[-1,1]^n}F=\min_{w\in\pmcube}F(w).
\end{equation}
The clauses define the objective only; clause satisfaction is not a
feasibility constraint.

\subsection{The node oracle and the general transfer}
A Boolean pseudoexpectation through degree $d$ is a linear functional
$\mathcal E$ on $\R[x]_{\le d}$ with $\mathcal E[1]=1$, respecting
$x_i^2=1$ through that degree and positive on $P^2$ whenever
$2\deg P\le d$. Boolean reduction replaces a monomial by the character
$\chi_A(x)=\prod_{i\in A}x_i$, where $A$ consists of its odd exponents;
$\chi_\varnothing=1$. Respecting the Boolean identities means that a
polynomial and its reduction have the same functional value.

For a nonempty coordinate product $B=\prod_i S_i\subseteq[-1,1]^n$, define
\begin{equation}\label{eq:xor-moment-hull}
 \mathcal Q(B)=\conv\{(x,xx^T):x\in B\}.
\end{equation}
Our order-$r$ node functional $L$ is defined through degree $2r$, normalized,
and satisfies
\begin{align}
 L\left[\left(\prod_{j=1}^a g_j\right)P^2\right]&\ge0,
 &\sum_j\deg g_j+2\deg P&\le2r,\label{eq:xor-node-preorder}\\
 L[qv]&=0,&\deg q+\deg v&\le2r,\label{eq:xor-node-local-equality}\\
 \bigl((L[x_i])_i,(L[x_ix_j])_{ij}\bigr)&\in\mathcal Q(B).
 &&\label{eq:xor-node-quadratic}
\end{align}
In \eqref{eq:xor-node-preorder}, each $g_j$ is any univariate polynomial
nonnegative on its coordinate set. Products may repeat factors, and the
empty product and globally coupled $P$ are allowed. In
\eqref{eq:xor-node-local-equality}, $q$ is any univariate polynomial
vanishing on its coordinate set, and $v$ is arbitrary through degree.
For boxes, all lower bounds below also cover the weaker oracle that uses
only the two bound slacks per coordinate, with arbitrary repeated products. Let $\operatorname{LB}_r(B)=\inf L[F]$, where $r\ge2$ makes the
cubic objective available. An empty feasible set of functionals has bound
$+\infty$.

For compact boxes, \eqref{eq:xor-node-quadratic} is equivalent to adding
\emph{linearly on moments} every quadratic polynomial inequality valid
on the box: separation applies to the compact convex hull in
\eqref{eq:xor-moment-hull}. Products or square localizers of arbitrary
coupled quadratic inequalities are not part of this grant. For nonclosed
coordinate sets we impose actual hull membership, which can be stronger
than all closed halfspace inequalities. No closure is taken.

\begin{theorem}[Quadratic-hull cover transfer]\label{thm:xor-transfer}
Let $r\ge2$ be an integer. Suppose a Boolean pseudoexpectation $\mathcal E$
through degree $4r$ satisfies
\begin{equation}\label{eq:xor-source-assumptions}
 \mathcal E[\chi_e]=b_e\quad(e\in E),\qquad
 \bigl(\mathcal E[x],\mathcal E[xx^T]\bigr)
 \in\conv\{(w,ww^T):w\in\pmcube\}.
\end{equation}
Every finite cover of $[-1,1]^n$ by coordinate products certified by
\eqref{eq:xor-node-preorder}--\eqref{eq:xor-node-quadratic} at a target
$T_*>0$ has at least
\begin{equation}\label{eq:xor-transfer-count}
 2^{mT_*/\Delta}
\end{equation}
regions. The coordinate sets may be disconnected, nonclosed, or have no
finite polynomial description.
\end{theorem}
\begin{proof}
Use all $2^n$ Boolean vertices as equally weighted witnesses. For a domain
$B$ containing a witness $w$, define
\[
 R=\{i:\{-1,1\}\not\subseteq S_i\},\qquad U=[n]\setminus R,
 \qquad L[p]=\mathcal E[p(w_R,x_U)].
\]
This fixes coordinates deterministically and then takes a marginal of
$\mathcal E$; it does not condition on an event under $\mathcal E$ and
needs no positive mass for the chosen signs. Substitution does not
increase degree, and $L[1]=1$.

First consider a box-slack product $gP^2$. Slacks in $R$ become
nonnegative constants. In $U$, the interval contains both endpoints and
therefore equals $[-1,1]$. Each remaining slack is $1+x_i$ or $1-x_i$.
Modulo Boolean identities, their product is zero if opposite signs occur
on a coordinate, and otherwise is a nonnegative multiple of an assignment
indicator $I$. Repeated factors change only the nonnegative multiplier.
If $a=\deg g$ and $d=\deg P$, the indicator has degree at most $a$.
Writing $\widehat P=P(w_R,x_U)$, Boolean idempotence gives
\begin{equation}\label{eq:xor-indicator-square}
 I\widehat P^2=(I\widehat P)^2\pmod{x_i^2-1},\qquad
 \deg(I\widehat P)\le a+d\le2r.
\end{equation}
Both sides are within degree $4r$, so Boolean reduction and square
positivity of $\mathcal E$ prove the inequality.

For the stronger local oracle, a factor on $R$ is again a nonnegative
constant. On $U$ its Boolean reduction is
\begin{equation}\label{eq:xor-local-endpoints}
 g(x_i)=g(1)\frac{1+x_i}{2}+g(-1)\frac{1-x_i}{2}
 \pmod{x_i^2-1},\qquad g(1),g(-1)\ge0.
\end{equation}
Remove constant factors and expand the product. Every summand is a
nonnegative multiple of a product of endpoint indicators. The number of
nonconstant factors, and hence the indicator degree, is at most the
original sum of generator degrees. Applying
\eqref{eq:xor-indicator-square} proves every allowed localizer, including
those with repeated factors. A local equality is zero
at the fixed witness on $R$, or has both endpoint values zero on $U$;
its Boolean reduction, also after any allowed multiplier, vanishes.

Let $\mu$ be a Boolean distribution realizing the original moments through
degree two. Marginalize $\mu$ on $U$ and set its other coordinates to
$w_R$. The resulting finite distribution is supported in $B$. Its means,
diagonal moments, and cross moments are exactly those of $L$, proving
\eqref{eq:xor-node-quadratic} even for nonclosed sets.

An untouched clause keeps cost zero. At most $\Delta|R|$ clauses touch
$R$, including multiplicity. After substitution each cost has the form
$(1-\sigma\chi_A)/2$ with $|A|\le3$. Since $\chi_A^2=1$,
\[
 0\le\mathcal E[(1\pm\chi_A)^2]
       =2(1\pm\mathcal E[\chi_A])
\]
proves its cost lies in $[0,1]$; these squares have degree at most six,
within $4r$. Thus
\begin{equation}\label{eq:xor-node-cost}
 \operatorname{LB}_r(B)\le L[F]\le\frac{\Delta|R|}{m}.
\end{equation}
Certification forces $|R|\ge mT_*/\Delta$. Every coordinate in $R$
permits exactly one Boolean endpoint when the domain contains a witness,
so the witness fraction in $B$ is $2^{-|R|}$ and is at most
$2^{-mT_*/\Delta}$. Domains containing no witnesses have fraction zero.
The union bound for any cover proves \eqref{eq:xor-transfer-count}.
\end{proof}

As in Section~\ref{subsec:certificate-model}, split points and branching
variables are arbitrary, and a complete split tree has at least this many
leaves. Regions discarded by objective propagation must also be certified
and charged; the theorem does not bound the leaves that remain after
uncertified deletion of feasible regions.

\subsection{Classical signed moments and an occurrence-controlled family}
\begin{lemma}[Realizing signed quadratic moments]\label{lem:signed-moment-law}
Let an augmented moment matrix indexed by $1,z_1,\ldots,z_N$ be positive
semidefinite, with all diagonal entries one and all entries in
$\{0,-1,1\}$. There is a probability distribution on $\{-1,1\}^N$
with exactly its prescribed first and second moments.
\end{lemma}
\begin{proof}
Represent the matrix as a Gram matrix of unit vectors $v_0,v_1,\ldots,v_N$.
Inner product $1$ or $-1$ forces equality or opposition. Partition the
vectors into classes up to sign and choose a representative for each.
Distinct representatives are orthogonal because their inner products
belong to $\{0,-1,1\}$ and cannot have absolute value one. Choose $v_0$
as representative of its own class. Give this class the deterministic
random value one; give every other class an independent unbiased sign.
If $v_j=\sigma_j v_{C(j)}$, set $Z_j=\sigma_j\xi_{C(j)}$.
A coordinate in the constant class has mean $\sigma_j$, and other
coordinates have mean zero. Two coordinates in one class have product
mean $\sigma_i\sigma_j$; different classes give product mean zero,
including when one is the constant class. These are exactly all Gram
entries, and the distribution has finite support.
\end{proof}

We state the source's width convention explicitly because the spatial
argument loses degree twice. The following lemma is the signed-character
construction in the proof of Schoenebeck's Lemma~13
\citep[Theorems 11--12 and Lemma 13]{Schoenebeck2008}, written as a
polynomial functional.

\begin{lemma}[Width and character degree]\label{lem:xor-width-moments}
Suppose a signed XOR formula has no width-$w$ XOR-resolution contradiction,
where $w\ge3$ is an integer and all initial clause widths are at most $w$.
There is a normalized Boolean functional through degree $w$ with every
character moment in $\{0,-1,1\}$, satisfying every clause in expectation
and positive on squares of degree at most $\lfloor w/2\rfloor$.
\end{lemma}
\begin{proof}
In multiplicative notation, the resolution rule derives
$\chi_{A\mathbin\triangle B}=st$ from $\chi_A=s$ and $\chi_B=t$
when $|A\mathbin\triangle B|\le w$. Let $\mathcal C_w$ be its closure,
including $\chi_\varnothing=1$. Absence of contradiction means that a
support can have at most one derived sign: opposite signs would derive
$\chi_\varnothing=-1$. Define $y_A$ for $|A|\le w$ to be that sign
when derived, and zero otherwise. In particular $y_\varnothing=1$.
Define $\mathcal E$ by Boolean reduction and $\mathcal E[\chi_A]=y_A$.
This proves normalization, signed moments, and the clause means.

Index characters by supports of size at most $\ell=\lfloor w/2\rfloor$.
Declare $I\sim J$ if a sign for $I\mathbin\triangle J$ is derived.
This is reflexive and symmetric. For transitivity, multiply the relations
for $I\mathbin\triangle J$ and $J\mathbin\triangle K$; the resulting
support $I\mathbin\triangle K$ has size at most $2\ell\le w$.
Thus the result is in the closure, with the product sign. Choose a
representative $I_0$ per class and orient $I$ by the derived sign
$s_I$ of $I\mathbin\triangle I_0$. Multiplication of two such relations
shows that $s_Is_J=y_{I\mathbin\triangle J}$ within a class. Give
distinct classes orthogonal unit vectors $e_C$ and put $v_I=s_Ie_{C(I)}$.
Then for all $|I|,|J|\le\ell$,
\begin{equation}\label{eq:xor-character-gram}
 \ip{v_I}{v_J}=y_{I\mathbin\triangle J}.
\end{equation}
For any polynomial $P$ of degree at most $\ell$, write its Boolean
reduction as $\sum_I p_I\chi_I$. Equation~\eqref{eq:xor-character-gram}
gives
\[
 \mathcal E[P^2]=\sum_{I,J}p_Ip_Jy_{I\mathbin\triangle J}
              =\left\|\sum_Ip_Iv_I\right\|^2\ge0.
\]
This proves positivity for arbitrary polynomial squares, including sums
whose terms jointly involve more than $\ell$ variables. The constant
character has a unit Gram vector, consistent with normalization.
\end{proof}

At density eight, Theorem~11 of \citet{Schoenebeck2008} gives a constant
$a>0$ such that width at least $an$ is available with probability $1-o(1)$,
after reducing $a$ to absorb integer rounding. Take the source parameters
$k=3$, $d=8$, $\delta=\gamma=1/4$, and its density exponent
$\varepsilon_{\rm src}=0$. Theorem~12's density condition is
$8\ge1+8\log2$, and the denominator $k/2-\gamma-1=1/4$ is positive.
Theorem~11 supplies the positive width constant; Lemma~13 converts width
$w$ to level $\lfloor w/2\rfloor$. Lemma~\ref{lem:xor-width-moments}
verifies directly that this means moments through degree $w$ and square
positivity through degree $\lfloor w/2\rfloor$, not moments only through
half the width. We need $4r\le w$ below and $4rD\le w$ for the lifts.
This choice avoids the boundary $\gamma=1/2$ in the source's illustrative
three-variable specialization.

\begin{theorem}[A bounded-occurrence cubic family]\label{thm:xor-family}
There are absolute constants $a>0$ and $n_0$ such that for every integer
$n\ge n_0$ there is an objective \eqref{eq:xor-objective} with
\begin{equation}\label{eq:xor-family-parameters}
 7n\le m\le8n,\qquad \Delta\le64,\qquad \OPT\ge1/8,
\end{equation}
and a Boolean functional of available degree at least $an$ as in
Lemma~\ref{lem:xor-width-moments}. For every integer $r\ge2$ with
$4r\le an$, any certificate at absolute tolerance $1/16$, or relative
tolerance $1/2$, under the quadratic-hull node oracle has at least
\begin{equation}\label{eq:xor-family-count}
 2^{7n/1024}
\end{equation}
regions. The same instance works simultaneously for all these orders.
\end{theorem}
\begin{proof}
Sample $m_0=8n$ independent clauses, each with a uniform three-element
support and an independent uniform sign. This is the random-XOR model in
Schoenebeck's definition: ordered distinct variables and independent literal
negations induce the same support and parity-sign distribution.
For a fixed Boolean assignment, its violation count $V$ is
$\operatorname{Bin}(8n,1/2)$, regardless of overlap among supports. For
$t>0$, the generating function and $\cosh u\le e^{u^2/2}$ give
\[
 \Prob\{V\le2n\}
 \le e^{-2nt}\E e^{-t(V-4n)}
 =e^{-2nt}\cosh(t/2)^{8n}
 \le e^{-2nt+nt^2}.
\]
The inequality for $\cosh$ follows by integrating
$(\log\cosh u)'=\tanh u\le u$ for $u\ge0$ and using symmetry.
Taking $t=1$ and a union bound over all assignments bounds the probability
of any assignment with at most $2n$ violations by
$2^ne^{-n}=e^{-(1-\log2)n}=o(1)$. Separately, the source gives the signed
functional with linear available degree with probability $1-o(1)$.

Let $D_i$ be the original number of clauses containing coordinate $i$;
$D_i\sim\operatorname{Bin}(8n,3/n)$. Delete every clause incident to a
coordinate with $D_i>64$. If $Z$ is the number deleted, then
$Z\le\sum_iD_i\ind_{D_i>64}$. Differentiating the binomial generating
polynomial yields
\begin{align}
 \E[D_i\ind_{D_i>64}]
 &\le2^{-64}\E[D_i2^{D_i}]
   =\frac{48(1+3/n)^{8n-1}}{2^{64}}\nonumber\\
 &\le\frac{48e^{24}}{2^{64}}<\frac18.
 \label{eq:xor-deletion}
\end{align}
For example $e<3$ makes the last inequality follow from
$384\cdot3^{24}<2^{64}$, an integer inequality. Hence
$\Prob\{Z>n\}\le\E Z/n<1/8$ by Markov's inequality.
No independence between different $D_i$ is used. For every sufficiently
large $n$, the sum of this failure bound and the two preceding $o(1)$
failure probabilities is less than one. Fix an instance in their
intersection.

At least $7n$ clauses remain, all remaining occurrences are at most $64$,
and every assignment still violates at least $n$ clauses. Equation
\eqref{eq:xor-vertex-min} gives $\OPT\ge n/m\ge1/8$.
Deleting clauses preserves the original functional, its positivity, all
character moments, and the means of retained clauses. Its degree-two
moments have an actual Boolean realization by
Lemma~\ref{lem:signed-moment-law}. Apply Theorem~\ref{thm:xor-transfer}
with any $4r\le an$. At the best incumbent the absolute target is
$\OPT-1/16\ge1/16$ and the relative target is $\OPT/2\ge1/16$.
Thus $mT_*/\Delta\ge7n/(16\cdot64)$, proving the count. Restricting
one source functional to lower degrees proves simultaneous validity.
\end{proof}

\begin{proposition}[Sensitivity and curvature]\label{prop:xor-sensitivity}
Throughout $[-1,1]^n$, the objective \eqref{eq:xor-objective} satisfies
\begin{equation}\label{eq:xor-sensitivity}
 |\partial_iF|\le\frac{\Delta}{2m},\qquad
 \sum_{j\ne i}|\partial_i\partial_jF|\le\frac{\Delta}{m},\qquad
 \|\nabla^2F\|_2\le\frac{\Delta}{m}.
\end{equation}
In the family of Theorem~\ref{thm:xor-family}, the last bound is at most
$64/(7n)$.
\end{proposition}
\begin{proof}
An incident clause contributes at most $1/(2m)$ to a first derivative
and at most two mixed second derivatives of that magnitude. The diagonal
second derivatives vanish. Summing absolute contributions gives the first
two bounds, and the symmetric Hessian's spectral norm is at most its
maximum absolute row sum.
\end{proof}
These estimates record the normalization of this family; they do not
assert that small curvature forces large certificates in general.
